% =====================================================================
%  CARNET D'ENTRAINEMENT — FICHE 6 — THÈME 2
%  Niveau FACILE — Exercices
% =====================================================================
\documentclass[11pt,a4paper]{article}
\usepackage[utf8]{inputenc}
\usepackage[T1]{fontenc}
\usepackage{lmodern}
\usepackage[french]{babel}
\usepackage{amsmath,amssymb,mathtools}
\usepackage{icomma}
\usepackage[version=4]{mhchem}
\usepackage{siunitx}
\sisetup{output-decimal-marker={,}}
\usepackage{xcolor}
\usepackage{colortbl}
\usepackage{tikz}
\usepackage[most]{tcolorbox}
\usepackage{enumitem}
\usepackage{array}
\usepackage{booktabs}
\usepackage{fancyhdr}
\usepackage[hidelinks]{hyperref}
\usetikzlibrary{arrows.meta,positioning,calc,shapes.geometric}
\usepackage[a4paper,margin=2.1cm,headheight=26pt,headsep=8pt]{geometry}

\definecolor{primary}{RGB}{146,95,160}
\definecolor{primarydark}{RGB}{92,52,104}
\definecolor{excol}{RGB}{0,135,90}

\tcbset{boxbase/.style={enhanced,breakable,boxrule=0.9pt,arc=2.5pt,
   left=3mm,right=3mm,top=2.2mm,bottom=2.2mm,
   fonttitle=\bfseries,coltitle=white,toptitle=0.6mm,bottomtitle=0.6mm}}
\newtcolorbox[auto counter]{exo}[1][]{boxbase,colback=primary!5,colframe=primary,
   title={\textsc{Exercice}~\thetcbcounter\ifx\relax#1\relax\relax\else\textmd{ —\itshape\ #1}\fi}}

% -------- macros des quatre grandeurs --------
\newcommand{\nmax}{n_{\max}(e^-)}
\newcommand{\nv}{n_{v}(e^-)}
\newcommand{\nl}{n_{\ell}(e^-)}
\newcommand{\nnl}{n_{n\ell}(e^-)}

\pagestyle{fancy}\fancyhf{}
\fancyhead[L]{\small\textcolor{primarydark}{\textbf{Fiche 6 · Thème 2} — Méthode des 5 étapes et structure de Lewis}}
\fancyhead[R]{\small\textcolor{primarydark}{\textsc{Facile}}}
\fancyfoot[C]{\small\thepage}
\renewcommand{\headrulewidth}{0.8pt}
\renewcommand{\headrule}{\hbox to\headwidth{\color{primary}\leaders\hrule height \headrulewidth\hfill}}
\setlength{\parindent}{0pt}\setlength{\parskip}{3pt}

\begin{document}
\begin{center}
{\Large\bfseries\color{primarydark}Thème 2 — Niveau Facile}\\[2pt]
{\small\itshape Lire la fiche \texttt{Tuto\_Theme-2} avant de commencer. Chaque exercice n'entraîne
qu'\emph{une} des quatre grandeurs à la fois. Aucune calculatrice n'est nécessaire.}
\end{center}
\medskip

\begin{exo}[calculer $\nmax = 8a + 2b$]
Pour chaque molécule, déterminer $a$ (atomes visant l'octet, soit tous sauf \ce{H}) et $b$
(atomes visant le duet, soit les \ce{H}), puis calculer $\nmax = 8a + 2b$.
\begin{enumerate}[label=\alph*),leftmargin=2em,itemsep=1pt]
  \item \ce{NH3} \qquad b) \ce{H2O} \qquad c) \ce{CH4} \qquad d) \ce{CO2}
\end{enumerate}
\end{exo}

\begin{exo}[calculer $\nv$ pour des composés neutres]
Les édifices ci-dessous sont \textbf{neutres} ($q=0$). Calculer $\nv = \sum N_v^{(i)}$ en
additionnant les électrons de valence de tous les atomes (\ce{H}{=}1, \ce{C}{=}4, \ce{N}{=}5,
\ce{O}{=}6, \ce{Cl}{=}7).
\begin{enumerate}[label=\alph*),leftmargin=2em,itemsep=1pt]
  \item \ce{NH3} \qquad b) \ce{H2O} \qquad c) \ce{CCl4} \qquad d) \ce{CO2}
\end{enumerate}
\end{exo}

\begin{exo}[du couple $(\nmax\,;\,\nv)$ au nombre de liaisons]
On \emph{donne} $\nmax$ et $\nv$. Calculer $\nl = \nmax - \nv$, puis le nombre de liaisons
$\nl/2$.
\begin{enumerate}[label=\alph*),leftmargin=2em,itemsep=1pt]
  \item \ce{NH3} : $(\nmax\,;\,\nv) = (14\,;\,8)$ ;
  \item \ce{CO2} : $(\nmax\,;\,\nv) = (24\,;\,16)$ ;
  \item \ce{H2O} : $(\nmax\,;\,\nv) = (12\,;\,8)$.
\end{enumerate}
\end{exo}

\begin{exo}[du couple $(\nv\,;\,\nl)$ au nombre de doublets libres]
On \emph{donne} $\nv$ et $\nl$. Calculer $\nnl = \nv - \nl$, puis le nombre de doublets non
liants $\nnl/2$.
\begin{enumerate}[label=\alph*),leftmargin=2em,itemsep=1pt]
  \item \ce{NH3} : $(\nv\,;\,\nl) = (8\,;\,6)$ ;
  \item \ce{CO2} : $(\nv\,;\,\nl) = (16\,;\,8)$ ;
  \item \ce{H2O} : $(\nv\,;\,\nl) = (8\,;\,4)$.
\end{enumerate}
\end{exo}

\begin{exo}[calculer $\nv$ pour des ions : gare au $-q$]
Ces édifices sont \textbf{chargés}. Calculer $\nv = \sum N_v^{(i)} - q$ en soignant le signe :
on \emph{ajoute} pour un anion, on \emph{retranche} pour un cation.
\begin{enumerate}[label=\alph*),leftmargin=2em,itemsep=1pt]
  \item \ce{OH^-} \qquad b) \ce{NH4+} \qquad c) \ce{CO3^2-} \qquad d) \ce{NO3^-}
\end{enumerate}
\end{exo}

\end{document}
